\documentclass[12pt]{report}
\usepackage[a4paper, margin=2.5cm]{geometry}
\pagenumbering{gobble}
\usepackage{hyperref}
\usepackage{tcolorbox}
\newcommand{\student}[3]{\begin{tcolorbox}\noindent Introduction to
    Psycholinguistics (WS2021-2022) \\ Report of \textbf{Paper~#3} by
    \textbf{#1 (#2)} 
  \end{tcolorbox}}



\begin{document}

%% Change:
\student{Shawon Ashraf}{3460184}{8}


\paragraph{APA reference:}
Federmeier, K. D., & Kutas, M. (1999). A Rose By Any Other name: Long-term Memory Structure and Sentence Processing. \textit{Journal of memory and Language}, 41, 469-495.


\paragraph{Research Question(s):}
\begin{enumerate}
    \item Is there an overlap between the effects of word specific contexts and semantic features while processing sentences?
    \item Does long term semantic memory structure play a role in this process?
\end{enumerate}
\vfill
\paragraph{Methodology:} 
The method involved 132 sentence pairs consisting of 3 target words at the end. These target words are of three types. The first is the expected target word, the second one being from the same semantic category as the expected word (within-category violation) and the third is a words which is semantically different (between-category violation) from the expected word. These target words are presented in a shuffled order each pair.  The first sentence in the pair works as a stimulus for the volunteers to infer the context and  the second sentence requires them to use the knowledge from the first one to complete it. Authors also recorded EEG of the volunteers while they were processing the sentences to measure brain activity. Alongside these procedures the authors used Cloze probability to measure the correctness of the predictions made by the volunteers. Also, they checked for the plausability of the selected target words if they belonged to the either of the violations. 
\vfill

\paragraph{Results:}
The main results state that when words belonging to both within-category and between-category violations have near or similar meaning, they have seemingly close responses to exhibit N400 effect. This indicates that in a long-term memory structure, the effects of semantic features and context can overlap. In other words, it can be said that when semantic knowledge is built up in the memory over a period of time, it takes precedence
\vfill

\paragraph{Personal Opinion:} 
The paper starts with the aim to show the how context and semantic feature can have an effect on sentence processing when set in a long-term memory setting. While the stimuli and the experiment measure that effect and the authors discuss on it as well, they focus the bulk of discussion on the EEG results and N400 effect, and do not manage to find concrete answers for some of the research questions. However the questions they bring in the discussion are worth exploring for future work, such as those on memory structure, overlaps and similarities between violations.
\vfill


\end{document}

